\section{Pipeline Story: From Measurement to Alignment}

\subsection{A Single Point}

\begin{remark}
Consider a single measured point \(x\) in the real world.
\end{remark}

\begin{remark}
It may originate from:
\begin{itemize}
\item a surveying campaign,
\item a track recording vehicle,
\item or a design dataset.
\end{itemize}
\end{remark}

\begin{remark}
At this stage, the point exists only in a global coordinate system.
\end{remark}

---

\subsection{Step 1: CRS Transformation}

\begin{remark}
The point is first interpreted within a coordinate reference system:
\[
x_{\mathrm{crs}}
\;\xrightarrow{\mathcal{C}}\;
x_{\mathrm{world}}.
\]
\end{remark}

\begin{remark}
This step ensures that all data is expressed in a consistent spatial
frame.
\end{remark}

---

\subsection{Step 2: Projection onto the Track}

\begin{remark}
The point is projected onto the alignment:
\[
x_{\mathrm{world}}
\;\xrightarrow{\mathcal{W}}\;
(s,d).
\]
\end{remark}

\begin{remark}
This is a nonlinear operation:
\begin{itemize}
\item the longitudinal position \(s\) is determined,
\item the lateral deviation \(d\) is computed.
\end{itemize}
\end{remark}

\begin{remark}
At this moment, the point becomes part of the intrinsic alignment
coordinate system.
\end{remark}

---

\subsection{Step 3: Interpretation through the Model}

\begin{remark}
The alignment model provides a curvature description:
\[
(s,d)
\;\xrightarrow{\mathcal{M}}\;
\kappa(s).
\]
\end{remark}

\begin{remark}
Here, two components interact:
\begin{itemize}
\item the parametric structure (design intent),
\item local corrections (data-driven adjustments).
\end{itemize}
\end{remark}

\begin{remark}
The point is no longer treated as isolated data, but as part of a
structured geometric system.
\end{remark}

---

\subsection{Step 4: Physical Realization}

\begin{remark}
The theoretical model is transformed by physical processes:
\[
\kappa(s)
\;\xrightarrow{\mathcal{R}}\;
\gamma_{\mathrm{real}}(s).
\]
\end{remark}

\begin{remark}
This step reflects:
\begin{itemize}
\item construction constraints,
\item mechanical behaviour,
\item maintenance operations.
\end{itemize}
\end{remark}

\begin{remark}
The resulting geometry is not identical to the theoretical model.
\end{remark}

---

\subsection{Step 5: Back to World Coordinates}

\begin{remark}
The realized alignment is mapped back into world space:
\[
\gamma_{\mathrm{real}}(s)
\;\xrightarrow{\mathcal{W}^{-1}}\;
x_{\mathrm{real}}.
\]
\end{remark}

\begin{remark}
This produces the geometry that is visible and measurable in reality.
\end{remark}

---

\subsection{Step 6: Comparison and Feedback}

\begin{remark}
The original measurement and the realized position are compared:
\[
x_{\mathrm{real}} \approx x_{\mathrm{world}}.
\]
\end{remark}

\begin{remark}
The deviation drives optimization:
\begin{itemize}
\item parameters are adjusted,
\item corrections are updated,
\item the pipeline is evaluated again.
\end{itemize}
\end{remark}

---

\subsection{Interpretation}

\begin{remark}
The point has passed through multiple representations:
\begin{itemize}
\item global coordinates,
\item intrinsic alignment coordinates,
\item parametric model,
\item realized geometry.
\end{itemize}
\end{remark}

\begin{remark}
Each transformation introduces structure, constraints, and
interpretation.
\end{remark}

---

\subsection{Key Insight}

\begin{proposition}
A measured point is not directly part of the alignment.

It becomes part of the alignment only after passing through the
world-to-track transformation and the alignment model.
\end{proposition}

---

\subsection{Connection to AIM}

\begin{summary}
The AIM framework transforms raw spatial data into structured alignment
information through a sequence of operators.

This process integrates:
\begin{itemize}
\item coordinate systems,
\item geometric modelling,
\item physical realization,
\item optimization.
\end{itemize}

Thus, alignment modelling is not a static representation, but a dynamic
process linking measurement, theory, and construction.
\end{summary}
